\cabstract{
    \par 量子纠缠现象具有深刻的物理意义，在高能对撞机条件下对粒子间的纠缠进行检验能够为理解这一现象的内涵提供崭新视角。量子层析技术通过构建微分散射截面与自旋关联矩阵之间的函数关系，已为实验上确定纠缠的存在与否提供了切实可行的计算工具。伴随着人工智能浪潮的兴起以及机器学习方法的推广应用，在粒子物理实验中检验粒子对的纠缠正逐渐成为可能。

    \par 本文主要利用基于扩散模型和Transformer模型构建的OmniLearn生成模型，对$pp \to \tau^+ \tau^- \to \left( \pi^+ \bar{\nu}_\tau)(\pi^- \nu_\tau \right)$蒙特卡洛模拟过程中末态粒子的运动学信息进行重建，并在完成参考系变换和对所有事件的统计后获得$\pi$介子动量方向与螺旋度基之间的夹角分布，进而计算出$\tau^+ \tau^-$二比特系综自旋关联矩阵的所有矩阵元。利用由Peres-Horodecki判据衍生出的纠缠判据$\mathcal{O}_E^\pm=0.1688>0$确定该过程中$\tau^+ \tau^-$粒子对的自旋之间存在纠缠。

    \par 研究中使用的机器学习方法，其所生成的各粒子运动学信息与蒙特卡洛模拟产生的重建级数据符合程度较高，相比传统的丢失质量计算器方法具备显著的优势。这一优势是本文自旋关联矩阵计算结果以及纠缠判断可信的有力佐证。

}{量子纠缠；对撞机；$\tau^+ \tau^-$；生成模型}
